\begin{figure}[H]
		\centering
		\tikzset{midarrow/.style={postaction={decorate},
				decoration={markings, mark=at position 0.58 with {\arrow{Stealth}}}}}
		\resizebox{0.66\textwidth}{!}{%
			\begin{tikzpicture}[>=Stealth, line width=0.8pt, scale=1.25,
				pt/.style={circle,fill,inner sep=1.3pt}]
				
				% proiezione obliqua:  (x,y,z) -> (x + 0.45 y , 0.26 y + z)
				\coordinate (O)    at (0,0);
				\coordinate (Ppi)  at (2.54,0.312);   % P_pi = proiezione di P nel piano
				\coordinate (P)    at (2.54,1.612);   % P sopra il piano
				\coordinate (Ad)   at (0,1.30);       % piede di d sull'asse (quota di P)
				\coordinate (C1) at (-1.95,-0.26);    % piano pi
				\coordinate (C2) at ( 3.55,-0.26);
				\coordinate (C3) at ( 4.90, 0.52);
				\coordinate (C4) at (-0.60, 0.52);
				
				% piano pi
				\draw[thin] (C1) -- (C2) -- (C3) -- (C4) -- cycle;
				\node at (-1.2,-0.06) {$\pi$};
				
				% asse per O parallelo a theta (tratteggiato) + theta assiale (doppia freccia)
				\draw[dashed, gray!60] (0,-0.7) -- (0,1.75);
				\draw[->>, thick] (0,1.75) -- (0,2.55) node[pos=0.6,left=2pt] {$\bm{\upvartheta}$};
				
				% componenti
				\draw[->, thick] (O) -- (Ppi) node[pos=0.80,below] {$\overrightarrow{OP}_{\!\perp}$};
				\draw[->, thick] (Ppi) -- (P) node[pos=0.5,right=1pt] {$\overrightarrow{OP}_{\!\parallel}$};
				
				% OP totale (freccia a meta', per non affollare P)
				\draw[thick, midarrow] (O) -- (P) node[pos=0.52,above left=-2pt] {$\overrightarrow{OP}$};
				
				% distanza d dall'asse (alla quota di P, parallela a OP_perp)
				\draw[thin] (Ad) -- (P) node[pos=0.5,above=1pt] {$d$};
				
				% angolo alpha in O fra OP e OP_perp
				\draw[->, thin] ($(O)+(7:0.7)$) arc[start angle=7, end angle=32.4, radius=0.7];
				\node at ($(O)+(19:0.92)$) {$\alpha$};
				
				% punti
				\node[pt,label={left:$O$}]           at (O)   {};
				\node[pt,label={right:$P$}]     at (P)   {};
				\node[pt,label={right:$P_\pi$}] at (Ppi) {};
				
			\end{tikzpicture}
		}
		\caption{Interpretazione geometrica del termine rotazionale
			$\bm{\upvartheta} \times \protect\overrightarrow{OP}$:
			la componente $\protect\overrightarrow{OP}_{\!\perp}$, ortogonale
			all'asse di rotazione, determina la distanza $d$ del punto $P$
			dall'asse; lo spostamento rotazionale ha modulo
			$|\bm{\upvartheta}|\,d$ ed è tangente alla circonferenza
			di raggio $d$.}
		\label{fig:termine_rotazionale}
	\end{figure}
